% Options for packages loaded elsewhere
\PassOptionsToPackage{unicode}{hyperref}
\PassOptionsToPackage{hyphens}{url}
\documentclass[
]{article}
\usepackage{xcolor}
\usepackage[margin=1in]{geometry}
\usepackage{amsmath,amssymb}
\setcounter{secnumdepth}{-\maxdimen} % remove section numbering
\usepackage{iftex}
\ifPDFTeX
  \usepackage[T1]{fontenc}
  \usepackage[utf8]{inputenc}
  \usepackage{textcomp} % provide euro and other symbols
\else % if luatex or xetex
  \usepackage{unicode-math} % this also loads fontspec
  \defaultfontfeatures{Scale=MatchLowercase}
  \defaultfontfeatures[\rmfamily]{Ligatures=TeX,Scale=1}
\fi
\usepackage{lmodern}
\ifPDFTeX\else
  % xetex/luatex font selection
\fi
% Use upquote if available, for straight quotes in verbatim environments
\IfFileExists{upquote.sty}{\usepackage{upquote}}{}
\IfFileExists{microtype.sty}{% use microtype if available
  \usepackage[]{microtype}
  \UseMicrotypeSet[protrusion]{basicmath} % disable protrusion for tt fonts
}{}
\makeatletter
\@ifundefined{KOMAClassName}{% if non-KOMA class
  \IfFileExists{parskip.sty}{%
    \usepackage{parskip}
  }{% else
    \setlength{\parindent}{0pt}
    \setlength{\parskip}{6pt plus 2pt minus 1pt}}
}{% if KOMA class
  \KOMAoptions{parskip=half}}
\makeatother
\usepackage{color}
\usepackage{fancyvrb}
\newcommand{\VerbBar}{|}
\newcommand{\VERB}{\Verb[commandchars=\\\{\}]}
\DefineVerbatimEnvironment{Highlighting}{Verbatim}{commandchars=\\\{\}}
% Add ',fontsize=\small' for more characters per line
\usepackage{framed}
\definecolor{shadecolor}{RGB}{248,248,248}
\newenvironment{Shaded}{\begin{snugshade}}{\end{snugshade}}
\newcommand{\AlertTok}[1]{\textcolor[rgb]{0.94,0.16,0.16}{#1}}
\newcommand{\AnnotationTok}[1]{\textcolor[rgb]{0.56,0.35,0.01}{\textbf{\textit{#1}}}}
\newcommand{\AttributeTok}[1]{\textcolor[rgb]{0.13,0.29,0.53}{#1}}
\newcommand{\BaseNTok}[1]{\textcolor[rgb]{0.00,0.00,0.81}{#1}}
\newcommand{\BuiltInTok}[1]{#1}
\newcommand{\CharTok}[1]{\textcolor[rgb]{0.31,0.60,0.02}{#1}}
\newcommand{\CommentTok}[1]{\textcolor[rgb]{0.56,0.35,0.01}{\textit{#1}}}
\newcommand{\CommentVarTok}[1]{\textcolor[rgb]{0.56,0.35,0.01}{\textbf{\textit{#1}}}}
\newcommand{\ConstantTok}[1]{\textcolor[rgb]{0.56,0.35,0.01}{#1}}
\newcommand{\ControlFlowTok}[1]{\textcolor[rgb]{0.13,0.29,0.53}{\textbf{#1}}}
\newcommand{\DataTypeTok}[1]{\textcolor[rgb]{0.13,0.29,0.53}{#1}}
\newcommand{\DecValTok}[1]{\textcolor[rgb]{0.00,0.00,0.81}{#1}}
\newcommand{\DocumentationTok}[1]{\textcolor[rgb]{0.56,0.35,0.01}{\textbf{\textit{#1}}}}
\newcommand{\ErrorTok}[1]{\textcolor[rgb]{0.64,0.00,0.00}{\textbf{#1}}}
\newcommand{\ExtensionTok}[1]{#1}
\newcommand{\FloatTok}[1]{\textcolor[rgb]{0.00,0.00,0.81}{#1}}
\newcommand{\FunctionTok}[1]{\textcolor[rgb]{0.13,0.29,0.53}{\textbf{#1}}}
\newcommand{\ImportTok}[1]{#1}
\newcommand{\InformationTok}[1]{\textcolor[rgb]{0.56,0.35,0.01}{\textbf{\textit{#1}}}}
\newcommand{\KeywordTok}[1]{\textcolor[rgb]{0.13,0.29,0.53}{\textbf{#1}}}
\newcommand{\NormalTok}[1]{#1}
\newcommand{\OperatorTok}[1]{\textcolor[rgb]{0.81,0.36,0.00}{\textbf{#1}}}
\newcommand{\OtherTok}[1]{\textcolor[rgb]{0.56,0.35,0.01}{#1}}
\newcommand{\PreprocessorTok}[1]{\textcolor[rgb]{0.56,0.35,0.01}{\textit{#1}}}
\newcommand{\RegionMarkerTok}[1]{#1}
\newcommand{\SpecialCharTok}[1]{\textcolor[rgb]{0.81,0.36,0.00}{\textbf{#1}}}
\newcommand{\SpecialStringTok}[1]{\textcolor[rgb]{0.31,0.60,0.02}{#1}}
\newcommand{\StringTok}[1]{\textcolor[rgb]{0.31,0.60,0.02}{#1}}
\newcommand{\VariableTok}[1]{\textcolor[rgb]{0.00,0.00,0.00}{#1}}
\newcommand{\VerbatimStringTok}[1]{\textcolor[rgb]{0.31,0.60,0.02}{#1}}
\newcommand{\WarningTok}[1]{\textcolor[rgb]{0.56,0.35,0.01}{\textbf{\textit{#1}}}}
\usepackage{graphicx}
\makeatletter
\newsavebox\pandoc@box
\newcommand*\pandocbounded[1]{% scales image to fit in text height/width
  \sbox\pandoc@box{#1}%
  \Gscale@div\@tempa{\textheight}{\dimexpr\ht\pandoc@box+\dp\pandoc@box\relax}%
  \Gscale@div\@tempb{\linewidth}{\wd\pandoc@box}%
  \ifdim\@tempb\p@<\@tempa\p@\let\@tempa\@tempb\fi% select the smaller of both
  \ifdim\@tempa\p@<\p@\scalebox{\@tempa}{\usebox\pandoc@box}%
  \else\usebox{\pandoc@box}%
  \fi%
}
% Set default figure placement to htbp
\def\fps@figure{htbp}
\makeatother
\setlength{\emergencystretch}{3em} % prevent overfull lines
\providecommand{\tightlist}{%
  \setlength{\itemsep}{0pt}\setlength{\parskip}{0pt}}
\usepackage{booktabs}
\usepackage{bookmark}
\IfFileExists{xurl.sty}{\usepackage{xurl}}{} % add URL line breaks if available
\urlstyle{same}
\hypersetup{
  pdftitle={Case Study 1},
  pdfauthor={Maximilian Kraus, Lukas Vater, Hannah Weichhart},
  hidelinks,
  pdfcreator={LaTeX via pandoc}}

\title{Case Study 1}
\usepackage{etoolbox}
\makeatletter
\providecommand{\subtitle}[1]{% add subtitle to \maketitle
  \apptocmd{\@title}{\par {\large #1 \par}}{}{}
}
\makeatother
\subtitle{AKSTA Statistical Computing}
\author{Maximilian Kraus, Lukas Vater, Hannah Weichhart}
\date{}

\begin{document}
\maketitle

\section{1. Ratio of Fibonacci numbers (1
point)}\label{ratio-of-fibonacci-numbers-1-point}

\subsection{a.}\label{a.}

Write two different R functions which return the sequence
\(r_{i}=F_{i+1}/F_i\) for \(i=1,\ldots,n\) where \(F_i\) is the \(i\)th
Fibonacci number, once using \texttt{for} and once using \texttt{while}.

\begin{Shaded}
\begin{Highlighting}[]
\CommentTok{\# for }
\NormalTok{fib\_for }\OtherTok{\textless{}{-}} \ControlFlowTok{function}\NormalTok{(n) \{}
\NormalTok{  f }\OtherTok{\textless{}{-}} \FunctionTok{c}\NormalTok{(}\DecValTok{1}\NormalTok{, }\DecValTok{1}\NormalTok{)}
  
  \ControlFlowTok{for}\NormalTok{ (i }\ControlFlowTok{in} \DecValTok{3}\SpecialCharTok{:}\NormalTok{(n }\SpecialCharTok{+} \DecValTok{1}\NormalTok{)) \{}
\NormalTok{    f[i] }\OtherTok{\textless{}{-}}\NormalTok{ f[i }\SpecialCharTok{{-}} \DecValTok{1}\NormalTok{] }\SpecialCharTok{+}\NormalTok{ f[i }\SpecialCharTok{{-}} \DecValTok{2}\NormalTok{]}
\NormalTok{  \}}
  
  \FunctionTok{return}\NormalTok{(f[}\DecValTok{2}\SpecialCharTok{:}\NormalTok{(n }\SpecialCharTok{+} \DecValTok{1}\NormalTok{)] }\SpecialCharTok{/}\NormalTok{ f[}\DecValTok{1}\SpecialCharTok{:}\NormalTok{n])}
\NormalTok{\}}
  
\CommentTok{\# while}
\NormalTok{fib\_while }\OtherTok{\textless{}{-}} \ControlFlowTok{function}\NormalTok{(n) \{}
\NormalTok{  f }\OtherTok{\textless{}{-}} \FunctionTok{c}\NormalTok{(}\DecValTok{1}\NormalTok{, }\DecValTok{1}\NormalTok{)}

  \ControlFlowTok{while}\NormalTok{ (}\FunctionTok{length}\NormalTok{(f) }\SpecialCharTok{\textless{}}\NormalTok{ n }\SpecialCharTok{+} \DecValTok{1}\NormalTok{) \{}
\NormalTok{    n\_jetzt }\OtherTok{\textless{}{-}} \FunctionTok{length}\NormalTok{(f)}
\NormalTok{    f[n\_jetzt }\SpecialCharTok{+} \DecValTok{1}\NormalTok{] }\OtherTok{\textless{}{-}}\NormalTok{ f[n\_jetzt] }\SpecialCharTok{+}\NormalTok{ f[n\_jetzt }\SpecialCharTok{{-}} \DecValTok{1}\NormalTok{]}
\NormalTok{  \}}
  
  \FunctionTok{return}\NormalTok{(f[}\DecValTok{2}\SpecialCharTok{:}\NormalTok{(n }\SpecialCharTok{+} \DecValTok{1}\NormalTok{)] }\SpecialCharTok{/}\NormalTok{ f[}\DecValTok{1}\SpecialCharTok{:}\NormalTok{n])}
\NormalTok{\}}
\end{Highlighting}
\end{Shaded}

\subsection{b.}\label{b.}

Benchmark the two functions for \(n=200\) and \(n=2000\) (you can use
package \textbf{microbenchmark} or package \textbf{bench} for this
purpose). Which function is faster?

\begin{Shaded}
\begin{Highlighting}[]
\NormalTok{bm\_200 }\OtherTok{\textless{}{-}} \FunctionTok{microbenchmark}\NormalTok{(}\FunctionTok{fib\_for}\NormalTok{(}\DecValTok{200}\NormalTok{), }\FunctionTok{fib\_while}\NormalTok{(}\DecValTok{200}\NormalTok{), }\AttributeTok{times =} \DecValTok{100}\NormalTok{)}
\NormalTok{bm\_2000 }\OtherTok{\textless{}{-}} \FunctionTok{microbenchmark}\NormalTok{(}\FunctionTok{fib\_for}\NormalTok{(}\DecValTok{2000}\NormalTok{), }\FunctionTok{fib\_while}\NormalTok{(}\DecValTok{2000}\NormalTok{), }\AttributeTok{times =} \DecValTok{100}\NormalTok{)}

\FunctionTok{print}\NormalTok{(bm\_200)}
\end{Highlighting}
\end{Shaded}

\begin{verbatim}
## Unit: microseconds
##            expr   min    lq    mean median     uq     max neval
##    fib_for(200)  64.4  77.1 156.637  86.80 112.65  5937.0   100
##  fib_while(200) 113.1 132.3 395.754 149.25 231.00 21942.6   100
\end{verbatim}

\begin{Shaded}
\begin{Highlighting}[]
\FunctionTok{print}\NormalTok{(bm\_2000)}
\end{Highlighting}
\end{Shaded}

\begin{verbatim}
## Unit: microseconds
##             expr   min     lq     mean median     uq    max neval
##    fib_for(2000) 448.2 485.75  715.851  569.9  802.6 2855.6   100
##  fib_while(2000) 777.3 837.50 1152.849  907.0 1241.4 4250.9   100
\end{verbatim}

\subsection{c.}\label{c.}

Plot the sequence for \(n=100\). For which \(n\) does it start to
stabilize to a value? What number does the sequence converge to?

\begin{Shaded}
\begin{Highlighting}[]
\NormalTok{r }\OtherTok{\textless{}{-}} \FunctionTok{fib\_for}\NormalTok{(}\DecValTok{100}\NormalTok{)}
\FunctionTok{plot}\NormalTok{(r, }\AttributeTok{type =} \StringTok{"l"}\NormalTok{, }\AttributeTok{main =} \StringTok{"fib\_for"}\NormalTok{, }\AttributeTok{xlab =} \StringTok{"n"}\NormalTok{, }\AttributeTok{ylab =} \StringTok{"value"}\NormalTok{)}
\end{Highlighting}
\end{Shaded}

\pandocbounded{\includegraphics[keepaspectratio]{CS1_files/figure-latex/unnamed-chunk-3-1.pdf}}

This sequence converges to the golden ratio (\textasciitilde1.618). It
stabilizes at \(n\) = 10.

\section{2. Gamma function (1 point)}\label{gamma-function-1-point}

\subsection{a.}\label{a.-1}

Write a function to compute the following for \(n\) a positive integer
using the \texttt{gamma} function in base R. \[
\rho_n =\frac{\Gamma((n-1)/2)}{\Gamma(1/2)\Gamma((n - 2)/2)}
\]

\begin{Shaded}
\begin{Highlighting}[]
\NormalTok{roh\_n }\OtherTok{\textless{}{-}} \ControlFlowTok{function}\NormalTok{(n)\{}
  \ControlFlowTok{if}\NormalTok{((}\FunctionTok{is.integer}\NormalTok{(n) }\SpecialCharTok{==} \ConstantTok{FALSE}\NormalTok{) }\SpecialCharTok{|}\NormalTok{ (n}\SpecialCharTok{\textless{}}\DecValTok{0}\NormalTok{))\{}
    \FunctionTok{print}\NormalTok{(}\StringTok{\textquotesingle{}n is not a positive integer!\textquotesingle{}}\NormalTok{)\}}
  \ControlFlowTok{else}\NormalTok{\{}
    \FunctionTok{gamma}\NormalTok{((n}\DecValTok{{-}1}\NormalTok{)}\SpecialCharTok{/}\DecValTok{2}\NormalTok{)}\SpecialCharTok{/}\NormalTok{(}\FunctionTok{gamma}\NormalTok{(}\DecValTok{1}\SpecialCharTok{/}\DecValTok{2}\NormalTok{)}\SpecialCharTok{*}\FunctionTok{gamma}\NormalTok{((n}\DecValTok{{-}2}\NormalTok{)}\SpecialCharTok{/}\DecValTok{2}\NormalTok{))\}}
\NormalTok{\}}
\end{Highlighting}
\end{Shaded}

\subsection{b.}\label{b.-1}

Try \(n = 2000\). What do you observe? Why do you think the observed
behavior happens?

\begin{Shaded}
\begin{Highlighting}[]
\FunctionTok{roh\_n}\NormalTok{(}\DecValTok{2000}\DataTypeTok{L}\NormalTok{)}
\end{Highlighting}
\end{Shaded}

\begin{verbatim}
## [1] NaN
\end{verbatim}

For large n the function returns a NaN value. This could be, because the
gamma function returns inf values for those n, as the resulting number
is too large for R to process.

\subsection{c.}\label{c.-1}

Write an implementation which can also deal with large values of
\(n>1000\).

\begin{Shaded}
\begin{Highlighting}[]
\NormalTok{lroh\_n }\OtherTok{\textless{}{-}} \ControlFlowTok{function}\NormalTok{(n)\{}
  \ControlFlowTok{if}\NormalTok{((}\FunctionTok{is.integer}\NormalTok{(n) }\SpecialCharTok{==} \ConstantTok{FALSE}\NormalTok{) }\SpecialCharTok{|}\NormalTok{ (n}\SpecialCharTok{\textless{}}\DecValTok{0}\NormalTok{))\{}
    \FunctionTok{print}\NormalTok{(}\StringTok{\textquotesingle{}n not a positive integer!\textquotesingle{}}\NormalTok{)\}}
  \ControlFlowTok{else}\NormalTok{\{}
    \FunctionTok{exp}\NormalTok{(}\FunctionTok{lgamma}\NormalTok{((n}\DecValTok{{-}1}\NormalTok{)}\SpecialCharTok{/}\DecValTok{2}\NormalTok{)}\SpecialCharTok{{-}}\NormalTok{(}\FunctionTok{lgamma}\NormalTok{(}\FloatTok{0.5}\NormalTok{)}\SpecialCharTok{+}\FunctionTok{lgamma}\NormalTok{((n}\DecValTok{{-}2}\NormalTok{)}\SpecialCharTok{/}\DecValTok{2}\NormalTok{)))\}}
\NormalTok{\}}
\FunctionTok{lroh\_n}\NormalTok{(}\DecValTok{2000}\DataTypeTok{L}\NormalTok{)}
\end{Highlighting}
\end{Shaded}

\begin{verbatim}
## [1] 17.83009
\end{verbatim}

Instead of the gamma function, here the lgamma function (the natural
logarithm of gamma) is used for the calculations and afterwards the
exponential function calculated. As lgamma results in lower values, R
can handle these calculations better.

\subsection{d.}\label{d.}

Plot \(\rho_n/\sqrt{n}\) for different values of \(n\). Can you guess
the limit of \(\rho_n/\sqrt{n}\) for \(n \rightarrow\infty\)?

\pandocbounded{\includegraphics[keepaspectratio]{CS1_files/figure-latex/unnamed-chunk-7-1.pdf}}
The limit of \(\rho_n/\sqrt{n}\) for \(n \rightarrow\infty\) seems to be
0.4.

\section{3. Weak law of large numbers (1.5
points)}\label{weak-law-of-large-numbers-1.5-points}

The law of large numbers states that if you repeat an experiment
independently a large number of times and average the result, what you
obtain should be close to the expected value. Formally, the weak law of
large numbers states that if \(X_1, X_2, \ldots X_n\) are independently
and identically distributed (iid) random variables with a finite
expected value \(E(X_i)=\mu<\infty\), then, for any \(\epsilon > 0\)
\(\lim_{n\rightarrow \infty} P(|\bar X_n-\mu|>\epsilon)=0\) where
\(\bar X_n=\frac{X_1+X_2+\ldots+X_n}{n}\) is the sample mean (the sample
mean converges in probability to the expected value).

\subsection{a.}\label{a.-2}

\begin{itemize}
\item
  In R, simulate rolling a fair 6-sided die \(n=10\,000\) times using
  vectorized functions.
\item
  After each roll, compute the cumulative mean of all previous rolls.
\item
  Plot the running average against the number of rolls and comment on
  the result.
\end{itemize}

\subsection{b.}\label{b.-2}

\begin{itemize}
\item
  Replicate the die-rolling experiment \(M=50\) times using a for loop.
  This will result in \(50\) paths which you can plot against the number
  of rolls.
\item
  Plot all 50 paths on the same graph to observe the variation in the
  running averages.
\end{itemize}

\subsection{c.}\label{c.-2}

What is the theoretical expected value of a fair 6-sided die?

\subsection{d.}\label{d.-1}

Comment on how the running mean of the paths behaves relative to the
expected value as \(n\) increases.

\subsection{e.}\label{e.}

\begin{itemize}
\item
  Modify the code in b. to use a while loop.
\item
  Stop the loop when the proportion of replications where the difference
  between the running average and the expected value lies outside the
  interval \([-\epsilon,\epsilon]\) is less than 0.05. Use
  \(\epsilon = 0.01\).
\item
  How high is \(n\) at the end of the loop?
\end{itemize}

\subsection{f.}\label{f.}

\begin{itemize}
\item
  Repeat the analysis in b, but instead of simulating from a discrete
  uniform distribution, simulate from a continuous standard cauchy
  distribution using \texttt{rcauchy()}. The (standard) cauchy
  distribution arises from the ratio of two independent mean zero
  (standard) normally distributed random variables.
\item
  Plot the running mean against increasing \(n\). Comment on whether the
  weak Law of Large Numbers seems to hold for this distribution. Explain
  why or why not.
\end{itemize}

\section{4. Central limit theorem (1.5
points)}\label{central-limit-theorem-1.5-points}

The central limit theorem states that, under certain conditions, the sum
of a large number of random variables is approximately normal. Consider
the case where \(X_1, X_2,\ldots,X_n\) are i.i.d. random variables with
finite expected values (\(E(X_i)=\mu<\infty\)) and variances
\((Var(X_i)=\sigma^2<\infty)\). The central limit theorem states that
the cumulative distribution function (CDF) of the normalized random
variable \[
Z_n = \frac{\bar X_n - \mu}{\sigma/\sqrt{n}}
\] converges to the standard normal CDF as \(n\rightarrow\infty\).

\subsection{a.}\label{a.-3}

Derive the expected value of the sample mean
\(\bar X_n=\frac{X_1+X_2+\ldots+X_n}{n}\).

\subsection{b.}\label{b.-3}

Derive the variance of the sample mean
\(\bar X_n=\frac{X_1+X_2+\ldots+X_n}{n}\).

\subsection{c.}\label{c.-3}

The central limit theorem makes few assumption on the distribution of
the random variables \(X_i\). Consider again the die-rolling experiment
where \(X_i\) is the random variable giving the number in roll \(i\).

\begin{itemize}
\item
  Simulate Die Rolls: For each value of \(n\) (where
  \(n=\{100,1000,10000\}\)), simulate rolling a die \(n\) times. Record
  the mean of the \(n\) rolls.
\item
  Repeat the above process \(M=100\) times for each value of \(n\). This
  will give you 100 sample means for each \(n\).
\item
  Plot a histogram of the 100 sample means for each value of \(n\).
  Observe the shape of the histograms.
\item
  Discuss how the distribution of the sample means changes as \(n\)
  increases. Comment on whether the histograms appear to approach a
  normal distribution as \(n\), illustrating the Central Limit Theorem.
\end{itemize}

\subsection{d.}\label{d.-2}

Compute expected value \(\mu\) and the variance \(\sigma^2\) for the
random variable \(X_i\) (giving the number in die-roll \(i\)), which
follows a discrete uniform distribution.

\subsection{e.}\label{e.-1}

Now that you computed \(\mu\) and \(\sigma^2\), consider \(Z_n\).

\begin{itemize}
\item
  For each value of \(n\) (where \(n=\{100,1000,10000\}\)) and for
  \(M=\{100,1000,10000\}\) replications, simulate the rolling of \(n\)
  dice and calculate the \(Z_n\).
\item
  Plot histograms of \(Z_n\) for each combination of \(n\) and \(M\).
\item
  Discuss how the distribution of \(Z_n\) changes as \(n\) and \(M\)
  increase. Comment on whether the distributions of \(Z_n\) appear to
  approach a normal distribution, illustrating the Central Limit
  Theorem.
\end{itemize}

\section{5. Readable and efficient code (1
points)}\label{readable-and-efficient-code-1-points}

Read over the code below.

\subsection{a.}\label{a.-4}

Explain (in text) what the code does.

The first Code fits the random, normally distributed vectors x and z by
a linear model, if enough observations fulfill the required size. The x
vector is then calculated with a sinusoidal function of the residuals.

The second Code performs a kind of cross validation, as one fourth of
the data is first omitted for the calculation of a linear model of two
random, normally distributed vectors x and y, and the predicted. For
each of the predictions the RMSE is calculated.

\subsection{b.}\label{b.-4}

Explain (in text) what you would change to make the code more readable.

A for loop shortens the code significantly and one could easily expand
the number of steps by increasing the range iterated over.

A for loop also shortens the code and by adding the variables nb\_sets
and stepsize the amount of cross-validations performed can also be
easily adapted.

\subsection{c.}\label{c.-4}

Change the code according to b. and wrap it in a function. This function
should have at most 10 lines (without adding commands to more lines such
as \texttt{x\ \textless{}-\ 1;\ y\ \textless{}-\ 2}. Such commands will
count as 2 lines!). Check that the function called on the same input
outputs the same as the provided code.

\begin{Shaded}
\begin{Highlighting}[]
\FunctionTok{set.seed}\NormalTok{(}\DecValTok{1}\NormalTok{)}
\NormalTok{x }\OtherTok{\textless{}{-}} \FunctionTok{rnorm}\NormalTok{(}\DecValTok{100}\NormalTok{)}
\NormalTok{z }\OtherTok{\textless{}{-}} \FunctionTok{rnorm}\NormalTok{(}\DecValTok{100}\NormalTok{)}
\CommentTok{\#Adapted Code 1}
\ControlFlowTok{for}\NormalTok{ (i }\ControlFlowTok{in} \DecValTok{1}\SpecialCharTok{:}\DecValTok{4}\NormalTok{)\{}
  \ControlFlowTok{if}\NormalTok{(}\FunctionTok{sum}\NormalTok{(x }\SpecialCharTok{\textgreater{}=} \FloatTok{0.001}\SpecialCharTok{*}\NormalTok{i)}\SpecialCharTok{\textless{}}\NormalTok{i)\{}
    \FunctionTok{stop}\NormalTok{(}\FunctionTok{paste}\NormalTok{(}\StringTok{"step"}\NormalTok{, i,}\StringTok{"requires"}\NormalTok{,i,}\StringTok{"observation(s) with value \textgreater{}="}\NormalTok{, }\FloatTok{0.001}\SpecialCharTok{*}\NormalTok{i))}
\NormalTok{  \}}
\NormalTok{  fit }\OtherTok{\textless{}{-}} \FunctionTok{lm}\NormalTok{(x }\SpecialCharTok{\textasciitilde{}}\NormalTok{ z)}
\NormalTok{  r }\OtherTok{\textless{}{-}}\NormalTok{ fit}\SpecialCharTok{$}\NormalTok{residuals}
\NormalTok{  x }\OtherTok{\textless{}{-}}\NormalTok{ i}\SpecialCharTok{*}\FunctionTok{sin}\NormalTok{(r) }\SpecialCharTok{+}\NormalTok{ i}\SpecialCharTok{*}\FloatTok{0.01}
\NormalTok{\}}
\NormalTok{x}
\end{Highlighting}
\end{Shaded}

\begin{verbatim}
##          1          2          3          4          5          6          7 
## -1.0992886  1.4796352 -0.6516207  0.7865093  3.7873388 -0.2055179  3.8327000 
##          8          9         10         11         12         13         14 
##  2.0871822  3.2219962 -3.0933769  0.6399144  4.0325676 -0.7808135 -0.8326365 
##         15         16         17         18         19         20         21 
##  0.7073736 -3.2414091 -2.8376440  0.9776866  1.5608532  3.0227535  1.0117977 
##         22         23         24         25         26         27         28 
##  1.8922308 -1.0474880 -0.4636773  2.8222909 -3.4535959 -3.9591095 -0.5208395 
##         29         30         31         32         33         34         35 
## -1.9366900  4.0313764  0.7254867 -3.7758444  4.0170200 -3.2639844 -0.4572754 
##         36         37         38         39         40         41         42 
## -2.5234159 -2.5371644 -3.4025549  0.6499197  1.7909888 -3.9557224 -3.5524879 
##         43         44         45         46         47         48         49 
##  2.0066135  3.2915276 -0.9643167 -0.8508112  3.8888448  1.7701559 -3.8041387 
##         50         51         52         53         54         55         56 
##  0.9535984  4.0340296 -1.0616984  3.8507707 -0.6058260  0.4776386  0.5361922 
##         57         58         59         60         61         62         63 
## -2.6133980 -0.5347483  3.0972939 -3.9543801  1.0621806 -3.1799992  2.4340472 
##         64         65         66         67         68         69         70 
## -2.1800910 -0.7436082  1.3185751 -0.4879213  0.4959666  0.8098709  0.7497175 
##         71         72         73         74         75         76         77 
##  3.9472930 -0.6991699  2.9570004 -0.4460233 -0.5400743  3.3793674 -1.9905970 
##         78         79         80         81         82         83         84 
## -2.7888504 -1.2126583 -0.9658716 -1.4647239 -3.9458928  0.7544273 -0.7574175 
##         85         86         87         88         89         90         91 
##  3.0949240  3.7880760  1.0486519 -3.3233027  3.9897095  3.1324172 -1.4436191 
##         92         93         94         95         96         97         98 
##  0.7777796  0.5976052  2.3300564  0.5549268  3.2229364 -0.2445148 -1.4071426 
##         99        100 
## -0.4052558 -1.9274082
\end{verbatim}

\begin{Shaded}
\begin{Highlighting}[]
\FunctionTok{set.seed}\NormalTok{(}\DecValTok{1}\NormalTok{) }
\NormalTok{x }\OtherTok{\textless{}{-}} \FunctionTok{rnorm}\NormalTok{(}\DecValTok{1000}\NormalTok{)}
\NormalTok{y }\OtherTok{\textless{}{-}} \DecValTok{2} \SpecialCharTok{+}\NormalTok{ x }\SpecialCharTok{+} \FunctionTok{rnorm}\NormalTok{(}\DecValTok{1000}\NormalTok{)}
\NormalTok{df }\OtherTok{\textless{}{-}} \FunctionTok{data.frame}\NormalTok{(x, y)}
\CommentTok{\#adapted Code 2}
\NormalTok{nb\_sets }\OtherTok{\textless{}{-}} \DecValTok{4}
\NormalTok{stepsize }\OtherTok{\textless{}{-}} \FunctionTok{nrow}\NormalTok{(df)}\SpecialCharTok{/}\NormalTok{nb\_sets}
\ControlFlowTok{for}\NormalTok{ (j }\ControlFlowTok{in} \DecValTok{1}\SpecialCharTok{:}\NormalTok{nb\_sets)\{}
  \FunctionTok{cat}\NormalTok{(}\StringTok{"Step"}\NormalTok{, j, }\StringTok{"}\SpecialCharTok{\textbackslash{}n}\StringTok{"}\NormalTok{)}
\NormalTok{  section }\OtherTok{\textless{}{-}}\NormalTok{ (}\DecValTok{1}\SpecialCharTok{+}\NormalTok{stepsize}\SpecialCharTok{*}\NormalTok{(j}\DecValTok{{-}1}\NormalTok{))}\SpecialCharTok{:}\NormalTok{(stepsize}\SpecialCharTok{*}\NormalTok{j)}
\NormalTok{  fit }\OtherTok{\textless{}{-}} \FunctionTok{lm}\NormalTok{(y}\SpecialCharTok{\textasciitilde{}}\NormalTok{x, }\AttributeTok{data =}\NormalTok{ df[}\SpecialCharTok{{-}}\NormalTok{section,])}
\NormalTok{  p1 }\OtherTok{\textless{}{-}} \FunctionTok{predict}\NormalTok{(fit, }\AttributeTok{newdata =}\NormalTok{ df[section,])}
\NormalTok{  r }\OtherTok{\textless{}{-}} \FunctionTok{sqrt}\NormalTok{(}\FunctionTok{mean}\NormalTok{((p1 }\SpecialCharTok{{-}}\NormalTok{ df[section, }\StringTok{\textquotesingle{}y\textquotesingle{}}\NormalTok{])}\SpecialCharTok{\^{}}\DecValTok{2}\NormalTok{))}
  \FunctionTok{print}\NormalTok{(r)}
\NormalTok{\}}
\end{Highlighting}
\end{Shaded}

\begin{verbatim}
## Step 1 
## [1] 1.024996
## Step 2 
## [1] 0.9952113
## Step 3 
## [1] 1.06855
## Step 4 
## [1] 1.070726
\end{verbatim}

\begin{Shaded}
\begin{Highlighting}[]
\FunctionTok{set.seed}\NormalTok{(}\DecValTok{1}\NormalTok{)}
\NormalTok{x }\OtherTok{\textless{}{-}} \FunctionTok{rnorm}\NormalTok{(}\DecValTok{100}\NormalTok{)}
\NormalTok{z }\OtherTok{\textless{}{-}} \FunctionTok{rnorm}\NormalTok{(}\DecValTok{100}\NormalTok{)}
\ControlFlowTok{if}\NormalTok{ (}\FunctionTok{sum}\NormalTok{(x }\SpecialCharTok{\textgreater{}=}\NormalTok{ .}\DecValTok{001}\NormalTok{) }\SpecialCharTok{\textless{}} \DecValTok{1}\NormalTok{) \{}
  \FunctionTok{stop}\NormalTok{(}\StringTok{"step 1 requires 1 observation(s) with value \textgreater{}= .001"}\NormalTok{)}
\NormalTok{\}}
\NormalTok{fit }\OtherTok{\textless{}{-}} \FunctionTok{lm}\NormalTok{(x }\SpecialCharTok{\textasciitilde{}}\NormalTok{ z)}
\NormalTok{r }\OtherTok{\textless{}{-}}\NormalTok{ fit}\SpecialCharTok{$}\NormalTok{residuals}
\NormalTok{x }\OtherTok{\textless{}{-}} \FunctionTok{sin}\NormalTok{(r) }\SpecialCharTok{+}\NormalTok{ .}\DecValTok{01}
\ControlFlowTok{if}\NormalTok{ (}\FunctionTok{sum}\NormalTok{(x }\SpecialCharTok{\textgreater{}=}\NormalTok{ .}\DecValTok{002}\NormalTok{) }\SpecialCharTok{\textless{}} \DecValTok{2}\NormalTok{) \{}
  \FunctionTok{stop}\NormalTok{(}\StringTok{"step 2 requires 2 observation(s) with value \textgreater{}= .002"}\NormalTok{)}
\NormalTok{\}}
\NormalTok{fit }\OtherTok{\textless{}{-}} \FunctionTok{lm}\NormalTok{(x }\SpecialCharTok{\textasciitilde{}}\NormalTok{ z)}
\NormalTok{r }\OtherTok{\textless{}{-}}\NormalTok{ fit}\SpecialCharTok{$}\NormalTok{residuals}
\NormalTok{x }\OtherTok{\textless{}{-}} \DecValTok{2} \SpecialCharTok{*} \FunctionTok{sin}\NormalTok{(r) }\SpecialCharTok{+}\NormalTok{ .}\DecValTok{02}
\ControlFlowTok{if}\NormalTok{ (}\FunctionTok{sum}\NormalTok{(x }\SpecialCharTok{\textgreater{}=}\NormalTok{ .}\DecValTok{003}\NormalTok{) }\SpecialCharTok{\textless{}} \DecValTok{3}\NormalTok{) \{}
  \FunctionTok{stop}\NormalTok{(}\StringTok{"step 3 requires 3 observation(s) with value \textgreater{}= .003"}\NormalTok{)}
\NormalTok{\}}
\NormalTok{fit }\OtherTok{\textless{}{-}} \FunctionTok{lm}\NormalTok{(x }\SpecialCharTok{\textasciitilde{}}\NormalTok{ z)}
\NormalTok{r }\OtherTok{\textless{}{-}}\NormalTok{ fit}\SpecialCharTok{$}\NormalTok{residuals}
\NormalTok{x }\OtherTok{\textless{}{-}} \DecValTok{3} \SpecialCharTok{*} \FunctionTok{sin}\NormalTok{(r) }\SpecialCharTok{+}\NormalTok{ .}\DecValTok{03}
\ControlFlowTok{if}\NormalTok{ (}\FunctionTok{sum}\NormalTok{(x }\SpecialCharTok{\textgreater{}=}\NormalTok{ .}\DecValTok{004}\NormalTok{) }\SpecialCharTok{\textless{}} \DecValTok{4}\NormalTok{) \{}
  \FunctionTok{stop}\NormalTok{(}\StringTok{"step 4 requires 4 observation(s) with value \textgreater{}= .004"}\NormalTok{)}
\NormalTok{\}}
\NormalTok{fit }\OtherTok{\textless{}{-}} \FunctionTok{lm}\NormalTok{(x }\SpecialCharTok{\textasciitilde{}}\NormalTok{ z)}
\NormalTok{r }\OtherTok{\textless{}{-}}\NormalTok{ fit}\SpecialCharTok{$}\NormalTok{residuals}
\NormalTok{x }\OtherTok{\textless{}{-}} \DecValTok{4} \SpecialCharTok{*} \FunctionTok{sin}\NormalTok{(r) }\SpecialCharTok{+}\NormalTok{ .}\DecValTok{04}
\NormalTok{x}
\end{Highlighting}
\end{Shaded}

\begin{Shaded}
\begin{Highlighting}[]
\FunctionTok{set.seed}\NormalTok{(}\DecValTok{1}\NormalTok{) }
\NormalTok{x }\OtherTok{\textless{}{-}} \FunctionTok{rnorm}\NormalTok{(}\DecValTok{1000}\NormalTok{)}
\NormalTok{y }\OtherTok{\textless{}{-}} \DecValTok{2} \SpecialCharTok{+}\NormalTok{ x }\SpecialCharTok{+} \FunctionTok{rnorm}\NormalTok{(}\DecValTok{1000}\NormalTok{)}
\NormalTok{df }\OtherTok{\textless{}{-}} \FunctionTok{data.frame}\NormalTok{(x, y)}

\FunctionTok{cat}\NormalTok{(}\StringTok{"Step"}\NormalTok{, }\DecValTok{1}\NormalTok{, }\StringTok{"}\SpecialCharTok{\textbackslash{}n}\StringTok{"}\NormalTok{)}
\NormalTok{fit1 }\OtherTok{\textless{}{-}} \FunctionTok{lm}\NormalTok{(y }\SpecialCharTok{\textasciitilde{}}\NormalTok{ x, }\AttributeTok{data =}\NormalTok{ df[}\SpecialCharTok{{-}}\NormalTok{(}\DecValTok{1}\SpecialCharTok{:}\DecValTok{250}\NormalTok{),])}
\NormalTok{p1 }\OtherTok{\textless{}{-}} \FunctionTok{predict}\NormalTok{(fit1, }\AttributeTok{newdata =}\NormalTok{ df[(}\DecValTok{1}\SpecialCharTok{:}\DecValTok{250}\NormalTok{),])}
\NormalTok{r }\OtherTok{\textless{}{-}} \FunctionTok{sqrt}\NormalTok{(}\FunctionTok{mean}\NormalTok{((p1 }\SpecialCharTok{{-}}\NormalTok{ df[(}\DecValTok{1}\SpecialCharTok{:}\DecValTok{250}\NormalTok{),}\StringTok{"y"}\NormalTok{])}\SpecialCharTok{\^{}}\DecValTok{2}\NormalTok{))  }

\FunctionTok{cat}\NormalTok{(}\StringTok{"Step"}\NormalTok{, }\DecValTok{2}\NormalTok{, }\StringTok{"}\SpecialCharTok{\textbackslash{}n}\StringTok{"}\NormalTok{)}
\NormalTok{fit2 }\OtherTok{\textless{}{-}} \FunctionTok{lm}\NormalTok{(y }\SpecialCharTok{\textasciitilde{}}\NormalTok{ x, }\AttributeTok{data =}\NormalTok{ df[}\SpecialCharTok{{-}}\NormalTok{(}\DecValTok{251}\SpecialCharTok{:}\DecValTok{500}\NormalTok{),])}
\NormalTok{p2 }\OtherTok{\textless{}{-}} \FunctionTok{predict}\NormalTok{(fit2, }\AttributeTok{newdata =}\NormalTok{ df[(}\DecValTok{251}\SpecialCharTok{:}\DecValTok{500}\NormalTok{),])}
\NormalTok{r }\OtherTok{\textless{}{-}} \FunctionTok{c}\NormalTok{(r, }\FunctionTok{sqrt}\NormalTok{(}\FunctionTok{mean}\NormalTok{((p2 }\SpecialCharTok{{-}}\NormalTok{ df[(}\DecValTok{251}\SpecialCharTok{:}\DecValTok{500}\NormalTok{),}\StringTok{"y"}\NormalTok{])}\SpecialCharTok{\^{}}\DecValTok{2}\NormalTok{)))}

\FunctionTok{cat}\NormalTok{(}\StringTok{"Step"}\NormalTok{, }\DecValTok{3}\NormalTok{, }\StringTok{"}\SpecialCharTok{\textbackslash{}n}\StringTok{"}\NormalTok{)}
\NormalTok{fit3 }\OtherTok{\textless{}{-}} \FunctionTok{lm}\NormalTok{(y }\SpecialCharTok{\textasciitilde{}}\NormalTok{ x, }\AttributeTok{data =}\NormalTok{ df[}\SpecialCharTok{{-}}\NormalTok{(}\DecValTok{501}\SpecialCharTok{:}\DecValTok{750}\NormalTok{),])}
\NormalTok{p3 }\OtherTok{\textless{}{-}} \FunctionTok{predict}\NormalTok{(fit3, }\AttributeTok{newdata =}\NormalTok{ df[(}\DecValTok{501}\SpecialCharTok{:}\DecValTok{750}\NormalTok{),])}
\NormalTok{r }\OtherTok{\textless{}{-}} \FunctionTok{c}\NormalTok{(r, }\FunctionTok{sqrt}\NormalTok{(}\FunctionTok{mean}\NormalTok{((p3 }\SpecialCharTok{{-}}\NormalTok{ df[(}\DecValTok{501}\SpecialCharTok{:}\DecValTok{750}\NormalTok{),}\StringTok{"y"}\NormalTok{])}\SpecialCharTok{\^{}}\DecValTok{2}\NormalTok{)))}

\FunctionTok{cat}\NormalTok{(}\StringTok{"Step"}\NormalTok{, }\DecValTok{4}\NormalTok{, }\StringTok{"}\SpecialCharTok{\textbackslash{}n}\StringTok{"}\NormalTok{)}
\NormalTok{fit4 }\OtherTok{\textless{}{-}} \FunctionTok{lm}\NormalTok{(y }\SpecialCharTok{\textasciitilde{}}\NormalTok{ x, }\AttributeTok{data =}\NormalTok{ df[}\SpecialCharTok{{-}}\NormalTok{(}\DecValTok{751}\SpecialCharTok{:}\DecValTok{1000}\NormalTok{),])}
\NormalTok{p4 }\OtherTok{\textless{}{-}} \FunctionTok{predict}\NormalTok{(fit4, }\AttributeTok{newdata =}\NormalTok{ df[(}\DecValTok{751}\SpecialCharTok{:}\DecValTok{1000}\NormalTok{),])}
\NormalTok{r }\OtherTok{\textless{}{-}} \FunctionTok{c}\NormalTok{(r, }\FunctionTok{sqrt}\NormalTok{(}\FunctionTok{mean}\NormalTok{((p4 }\SpecialCharTok{{-}}\NormalTok{ df[(}\DecValTok{751}\SpecialCharTok{:}\DecValTok{1000}\NormalTok{),}\StringTok{"y"}\NormalTok{])}\SpecialCharTok{\^{}}\DecValTok{2}\NormalTok{)))  }
\NormalTok{r}
\end{Highlighting}
\end{Shaded}

\section{6. Measuring and improving performance (2
points)}\label{measuring-and-improving-performance-2-points}

Have a look at the code of the function below. It is a function for
performing a Kruskal-Wallis test, a robust non-parametric method for
testing whether samples come from the same distribution. (Note: we
assume no missing values are present in \texttt{x}).

\begin{Shaded}
\begin{Highlighting}[]
\NormalTok{kwtest }\OtherTok{\textless{}{-}} \ControlFlowTok{function}\NormalTok{ (x, g, ...) }
\NormalTok{\{}
    \ControlFlowTok{if}\NormalTok{ (}\FunctionTok{is.list}\NormalTok{(x)) \{}
        \ControlFlowTok{if}\NormalTok{ (}\FunctionTok{length}\NormalTok{(x) }\SpecialCharTok{\textless{}} \DecValTok{2}\DataTypeTok{L}\NormalTok{) }
            \FunctionTok{stop}\NormalTok{(}\StringTok{"\textquotesingle{}x\textquotesingle{} must be a list with at least 2 elements"}\NormalTok{)}
        \ControlFlowTok{if}\NormalTok{ (}\SpecialCharTok{!}\FunctionTok{missing}\NormalTok{(g)) }
            \FunctionTok{warning}\NormalTok{(}\StringTok{"\textquotesingle{}x\textquotesingle{} is a list, so ignoring argument \textquotesingle{}g\textquotesingle{}"}\NormalTok{)}
        \ControlFlowTok{if}\NormalTok{ (}\SpecialCharTok{!}\FunctionTok{all}\NormalTok{(}\FunctionTok{sapply}\NormalTok{(x, is.numeric))) }
            \FunctionTok{warning}\NormalTok{(}\StringTok{"some elements of \textquotesingle{}x\textquotesingle{} are not numeric and will be coerced to numeric"}\NormalTok{)}
\NormalTok{        k }\OtherTok{\textless{}{-}} \FunctionTok{length}\NormalTok{(x)}
\NormalTok{        l }\OtherTok{\textless{}{-}} \FunctionTok{lengths}\NormalTok{(x)}
        \ControlFlowTok{if}\NormalTok{ (}\FunctionTok{any}\NormalTok{(l }\SpecialCharTok{==} \DecValTok{0}\DataTypeTok{L}\NormalTok{)) }
            \FunctionTok{stop}\NormalTok{(}\StringTok{"all groups must contain data"}\NormalTok{)}
\NormalTok{        g }\OtherTok{\textless{}{-}} \FunctionTok{factor}\NormalTok{(}\FunctionTok{rep.int}\NormalTok{(}\FunctionTok{seq\_len}\NormalTok{(k), l))}
\NormalTok{        x }\OtherTok{\textless{}{-}} \FunctionTok{unlist}\NormalTok{(x)}
\NormalTok{    \}}
    \ControlFlowTok{else}\NormalTok{ \{}
        \ControlFlowTok{if}\NormalTok{ (}\FunctionTok{length}\NormalTok{(x) }\SpecialCharTok{!=} \FunctionTok{length}\NormalTok{(g)) }
            \FunctionTok{stop}\NormalTok{(}\StringTok{"\textquotesingle{}x\textquotesingle{} and \textquotesingle{}g\textquotesingle{} must have the same length"}\NormalTok{)}
\NormalTok{        g }\OtherTok{\textless{}{-}} \FunctionTok{factor}\NormalTok{(g)}
\NormalTok{        k }\OtherTok{\textless{}{-}} \FunctionTok{nlevels}\NormalTok{(g)}
        \ControlFlowTok{if}\NormalTok{ (k }\SpecialCharTok{\textless{}} \DecValTok{2}\DataTypeTok{L}\NormalTok{) }
            \FunctionTok{stop}\NormalTok{(}\StringTok{"all observations are in the same group"}\NormalTok{)}
\NormalTok{    \}}
\NormalTok{    n }\OtherTok{\textless{}{-}} \FunctionTok{length}\NormalTok{(x)}
    \ControlFlowTok{if}\NormalTok{ (n }\SpecialCharTok{\textless{}} \DecValTok{2}\DataTypeTok{L}\NormalTok{) }
        \FunctionTok{stop}\NormalTok{(}\StringTok{"not enough observations"}\NormalTok{)}
\NormalTok{    r }\OtherTok{\textless{}{-}} \FunctionTok{rank}\NormalTok{(x)}
\NormalTok{    TIES }\OtherTok{\textless{}{-}} \FunctionTok{table}\NormalTok{(x)}
\NormalTok{    STATISTIC }\OtherTok{\textless{}{-}} \FunctionTok{sum}\NormalTok{(}\FunctionTok{tapply}\NormalTok{(r, g, sum)}\SpecialCharTok{\^{}}\DecValTok{2}\SpecialCharTok{/}\FunctionTok{tapply}\NormalTok{(r, g, length))}
\NormalTok{    STATISTIC }\OtherTok{\textless{}{-}}\NormalTok{ ((}\DecValTok{12} \SpecialCharTok{*}\NormalTok{ STATISTIC}\SpecialCharTok{/}\NormalTok{(n }\SpecialCharTok{*}\NormalTok{ (n }\SpecialCharTok{+} \DecValTok{1}\NormalTok{)) }\SpecialCharTok{{-}} \DecValTok{3} \SpecialCharTok{*}\NormalTok{ (n }\SpecialCharTok{+} \DecValTok{1}\NormalTok{))}\SpecialCharTok{/}\NormalTok{(}\DecValTok{1} \SpecialCharTok{{-}} 
        \FunctionTok{sum}\NormalTok{(TIES}\SpecialCharTok{\^{}}\DecValTok{3} \SpecialCharTok{{-}}\NormalTok{ TIES)}\SpecialCharTok{/}\NormalTok{(n}\SpecialCharTok{\^{}}\DecValTok{3} \SpecialCharTok{{-}}\NormalTok{ n)))}
\NormalTok{    PARAMETER }\OtherTok{\textless{}{-}}\NormalTok{ k }\SpecialCharTok{{-}} \DecValTok{1}\DataTypeTok{L}
\NormalTok{    PVAL }\OtherTok{\textless{}{-}} \FunctionTok{pchisq}\NormalTok{(STATISTIC, PARAMETER, }\AttributeTok{lower.tail =} \ConstantTok{FALSE}\NormalTok{)}
    \FunctionTok{names}\NormalTok{(STATISTIC) }\OtherTok{\textless{}{-}} \StringTok{"Kruskal{-}Wallis chi{-}squared"}
    \FunctionTok{names}\NormalTok{(PARAMETER) }\OtherTok{\textless{}{-}} \StringTok{"df"}
\NormalTok{    RVAL }\OtherTok{\textless{}{-}} \FunctionTok{list}\NormalTok{(}\AttributeTok{statistic =}\NormalTok{ STATISTIC, }\AttributeTok{parameter =}\NormalTok{ PARAMETER, }
        \AttributeTok{p.value =}\NormalTok{ PVAL, }\AttributeTok{method =} \StringTok{"Kruskal{-}Wallis rank sum test"}\NormalTok{)}
    \FunctionTok{return}\NormalTok{(RVAL)}
\NormalTok{\}}
\end{Highlighting}
\end{Shaded}

\subsection{a.}\label{a.-5}

Write a pseudo code outlining what the function does.

\begin{enumerate}
\def\labelenumi{\arabic{enumi}.}
\item
  Input-Handling/Validation

  IF x is a list:

  \begin{itemize}
  \tightlist
  \item
    Check for having at least 2 elements and if all elements are numeric
  \item
    Ensure that no element is empty
  \item
    Create a grouping factor
  \item
    Flatten list into a single vector
  \end{itemize}

  ELSE IF x is a vector:

  \begin{itemize}
  \tightlist
  \item
    Ensure x and g have the same length
  \item
    Convert g into a factor
  \item
    Check for at least 2 groups
  \end{itemize}
\item
  Pre-processing
\end{enumerate}

\begin{itemize}
\tightlist
\item
  Confirm that total number fo observations is \textgreater=2
\item
  Assign ranks to all values
\item
  Identify ties
\end{itemize}

\begin{enumerate}
\def\labelenumi{\arabic{enumi}.}
\setcounter{enumi}{2}
\tightlist
\item
  Statistical calculation
\end{enumerate}

\begin{itemize}
\tightlist
\item
  Calculate the sum of ranks for each group
\item
  Compute the Kruskal-Wallis test statistic
\item
  Adjust the statistic
\end{itemize}

\begin{enumerate}
\def\labelenumi{\arabic{enumi}.}
\setcounter{enumi}{3}
\tightlist
\item
  Significance test
\end{enumerate}

\begin{itemize}
\tightlist
\item
  Determine degrees of freedom
\item
  Calculate p-value
\end{itemize}

\begin{enumerate}
\def\labelenumi{\arabic{enumi}.}
\setcounter{enumi}{4}
\tightlist
\item
  Output
\end{enumerate}

\begin{itemize}
\tightlist
\item
  Combine the statistic
\item
  Return list
\end{itemize}

\subsection{b.}\label{b.-5}

For example data, call the function in two ways: once using \texttt{x}
as a list and once using \texttt{x} as a vector with a corresponding
\texttt{g} argument. Ensure that the two different function calls return
the same thing by aligning the inputs.

\begin{Shaded}
\begin{Highlighting}[]
\CommentTok{\# input as list}
\NormalTok{x\_list }\OtherTok{\textless{}{-}} \FunctionTok{list}\NormalTok{(}
  \AttributeTok{group1 =} \FunctionTok{c}\NormalTok{(}\FloatTok{2.4}\NormalTok{, }\FloatTok{2.8}\NormalTok{, }\FloatTok{3.0}\NormalTok{),}
  \AttributeTok{group2 =} \FunctionTok{c}\NormalTok{(}\FloatTok{3.5}\NormalTok{, }\FloatTok{3.8}\NormalTok{, }\FloatTok{4.0}\NormalTok{),}
  \AttributeTok{group3 =} \FunctionTok{c}\NormalTok{(}\FloatTok{4.5}\NormalTok{, }\FloatTok{4.9}\NormalTok{, }\FloatTok{5.1}\NormalTok{)}
\NormalTok{)}

\CommentTok{\# input as vector}
\NormalTok{x\_vector }\OtherTok{=} \FunctionTok{c}\NormalTok{(}\FloatTok{2.4}\NormalTok{, }\FloatTok{2.8}\NormalTok{, }\FloatTok{3.0}\NormalTok{, }\FloatTok{3.5}\NormalTok{, }\FloatTok{3.8}\NormalTok{, }\FloatTok{4.0}\NormalTok{, }\FloatTok{4.5}\NormalTok{, }\FloatTok{4.9}\NormalTok{, }\FloatTok{5.1}\NormalTok{)}
\NormalTok{g\_vector }\OtherTok{=} \FunctionTok{factor}\NormalTok{(}\FunctionTok{c}\NormalTok{(}\StringTok{"group1"}\NormalTok{, }\StringTok{"group1"}\NormalTok{, }\StringTok{"group1"}\NormalTok{,}
                    \StringTok{"group2"}\NormalTok{, }\StringTok{"group2"}\NormalTok{, }\StringTok{"group2"}\NormalTok{,}
                    \StringTok{"group3"}\NormalTok{, }\StringTok{"group3"}\NormalTok{, }\StringTok{"group3"}\NormalTok{))}

\NormalTok{result\_list }\OtherTok{\textless{}{-}} \FunctionTok{kwtest}\NormalTok{(x\_list)}
\NormalTok{result\_vector }\OtherTok{\textless{}{-}} \FunctionTok{kwtest}\NormalTok{(x\_vector, g\_vector)}

\CommentTok{\# check if the outputs are equal}
\FunctionTok{all.equal}\NormalTok{(result\_list, result\_vector)}
\end{Highlighting}
\end{Shaded}

\begin{verbatim}
## [1] TRUE
\end{verbatim}

\subsection{c.}\label{c.-5}

Make a faster version of \texttt{kwtest()} that only computes the
Kruskal-Wallis \textbf{test statistic} when the input is a numeric
variable \texttt{x} and a variable \texttt{g} which gives the group
membership. You can try simplifying the function above or by coding from
the mathematical definition (see
\url{https://en.wikipedia.org/wiki/Kruskal\%E2\%80\%93Wallis_one-way_analysis_of_variance}).
This function should also perform some checks to ensure the correctness
of the inputs (use \texttt{kwtest()} as inspiration).

\begin{Shaded}
\begin{Highlighting}[]
\NormalTok{kwfast }\OtherTok{\textless{}{-}} \ControlFlowTok{function}\NormalTok{(x, g) \{}
  
  \CommentTok{\# basic checks}
  \ControlFlowTok{if}\NormalTok{ (}\SpecialCharTok{!}\FunctionTok{is.numeric}\NormalTok{(x)) \{}
\NormalTok{    x }\OtherTok{\textless{}{-}} \FunctionTok{as.numeric}\NormalTok{(x)}
\NormalTok{  \}}
  
  \ControlFlowTok{if}\NormalTok{ (}\FunctionTok{length}\NormalTok{(x) }\SpecialCharTok{!=} \FunctionTok{length}\NormalTok{(g)) \{}
    \FunctionTok{stop}\NormalTok{(}\StringTok{"\textquotesingle{}x\textquotesingle{} and \textquotesingle{}g\textquotesingle{} must have the same length"}\NormalTok{)}
\NormalTok{  \}}
  
  \CommentTok{\# identify groups}
\NormalTok{  g }\OtherTok{\textless{}{-}} \FunctionTok{as.factor}\NormalTok{(g)}
\NormalTok{  n }\OtherTok{\textless{}{-}} \FunctionTok{length}\NormalTok{(x)}
\NormalTok{  k }\OtherTok{\textless{}{-}} \FunctionTok{nlevels}\NormalTok{(g)}
  
  \CommentTok{\# pre{-}processing}
\NormalTok{  r }\OtherTok{\textless{}{-}} \FunctionTok{rank}\NormalTok{(x)}
  
  \CommentTok{\# stat calculation}
\NormalTok{  rank\_sums }\OtherTok{\textless{}{-}} \FunctionTok{rowsum}\NormalTok{(r, g)}
\NormalTok{  group\_sizes }\OtherTok{\textless{}{-}} \FunctionTok{as.numeric}\NormalTok{(}\FunctionTok{table}\NormalTok{(g))}
  
  \CommentTok{\# kruskal{-}wallis formula}
\NormalTok{  statistic }\OtherTok{\textless{}{-}}\NormalTok{ (}\DecValTok{12} \SpecialCharTok{/}\NormalTok{ (n }\SpecialCharTok{*}\NormalTok{ (n }\SpecialCharTok{+} \DecValTok{1}\NormalTok{))) }\SpecialCharTok{*} \FunctionTok{sum}\NormalTok{(rank\_sums}\SpecialCharTok{\^{}}\DecValTok{2} \SpecialCharTok{/}\NormalTok{ group\_sizes) }\SpecialCharTok{{-}} \DecValTok{3} \SpecialCharTok{*}\NormalTok{ (n }\SpecialCharTok{+} \DecValTok{1}\NormalTok{)}
  
  \FunctionTok{return}\NormalTok{(statistic)}
\NormalTok{\}}
\end{Highlighting}
\end{Shaded}

\subsection{d.}\label{d.-3}

Consider the following scenario. You have samples available from
multiple experiments \(m=1000\) where you collect the numerical values
for the quantity of interest \texttt{x} and the group membership for
\(n\) individuals. The first \(20\) individuals in each sample belong to
group 1, the following \(20\) individuals in each sample belong to group
2, the last \(10\) individuals in each sample belong to group 3. Use the
following code to simulate such a data structure:

\begin{Shaded}
\begin{Highlighting}[]
\FunctionTok{set.seed}\NormalTok{(}\DecValTok{1234}\NormalTok{)}
\NormalTok{m }\OtherTok{\textless{}{-}} \DecValTok{1000} \CommentTok{\# number of repetitions}
\NormalTok{n }\OtherTok{\textless{}{-}} \DecValTok{50}   \CommentTok{\# number of individuals}
\NormalTok{X }\OtherTok{\textless{}{-}} \FunctionTok{matrix}\NormalTok{(}\FunctionTok{rt}\NormalTok{(m }\SpecialCharTok{*}\NormalTok{ n, }\AttributeTok{df =} \DecValTok{10}\NormalTok{), }\AttributeTok{nrow =}\NormalTok{ m)}
\NormalTok{grp }\OtherTok{\textless{}{-}} \FunctionTok{rep}\NormalTok{(}\DecValTok{1}\SpecialCharTok{:}\DecValTok{3}\NormalTok{, }\FunctionTok{c}\NormalTok{(}\DecValTok{20}\NormalTok{, }\DecValTok{20}\NormalTok{, }\DecValTok{10}\NormalTok{))}
\end{Highlighting}
\end{Shaded}

Write a function which performs the Kruskal-Wallis test using the
function \texttt{stats:::kruskal.test.default()} for \(m\) repeated
experiments and returns a vector of \(m\) test statistics. The input of
this function are a matrix \texttt{X} with \(m\) rows and \(n\) columns
and a vector \texttt{g} which gives the grouping.

\begin{Shaded}
\begin{Highlighting}[]
\NormalTok{run\_kruskal\_default }\OtherTok{\textless{}{-}} \ControlFlowTok{function}\NormalTok{(X, g) \{}
  
\NormalTok{  results }\OtherTok{\textless{}{-}} \FunctionTok{apply}\NormalTok{(X, }\DecValTok{1}\NormalTok{, }\ControlFlowTok{function}\NormalTok{(row) \{}
    
\NormalTok{    test\_obj }\OtherTok{\textless{}{-}}\NormalTok{ stats}\SpecialCharTok{:::}\FunctionTok{kruskal.test.default}\NormalTok{(row, g)}
    
    \FunctionTok{return}\NormalTok{(test\_obj}\SpecialCharTok{$}\NormalTok{statistic)}
\NormalTok{  \})}
  
  \FunctionTok{return}\NormalTok{(}\FunctionTok{as.numeric}\NormalTok{(results))}
\NormalTok{\}}
\end{Highlighting}
\end{Shaded}

\subsection{e.}\label{e.-2}

Write a function which performs the Kruskal-Wallis test using the
function in point c.~for \(m\) repeated experiments and returns a vector
of \(m\) test statistics. The input of this function are a matrix
\texttt{X} with \(m\) rows and \(n\) columns and a vector \texttt{g}
which gives the grouping.

\begin{Shaded}
\begin{Highlighting}[]
\NormalTok{run\_kruskal\_fast }\OtherTok{\textless{}{-}} \ControlFlowTok{function}\NormalTok{(X, g) \{}
  
\NormalTok{  results }\OtherTok{\textless{}{-}} \FunctionTok{apply}\NormalTok{(X, }\DecValTok{1}\NormalTok{, }\ControlFlowTok{function}\NormalTok{(row) \{}
    
    \FunctionTok{kwfast}\NormalTok{(}\AttributeTok{x =}\NormalTok{ row, }\AttributeTok{g =}\NormalTok{ g)}
\NormalTok{  \})}
  
  \FunctionTok{return}\NormalTok{(results)}
\NormalTok{\}}
\end{Highlighting}
\end{Shaded}

\subsection{f.}\label{f.-1}

Compare the performance of the two approaches using a benchmarking
package on the data generated above. Comment on the results.

\begin{Shaded}
\begin{Highlighting}[]
\NormalTok{performance\_comparison }\OtherTok{\textless{}{-}} \FunctionTok{microbenchmark}\NormalTok{(}
  \AttributeTok{default\_method =} \FunctionTok{run\_kruskal\_default}\NormalTok{(X, grp),}
  \AttributeTok{fast\_method    =} \FunctionTok{run\_kruskal\_fast}\NormalTok{(X, grp),}
  \AttributeTok{times =} \DecValTok{100}
\NormalTok{)}

\FunctionTok{print}\NormalTok{(performance\_comparison)}
\end{Highlighting}
\end{Shaded}

\begin{verbatim}
## Unit: milliseconds
##            expr      min       lq     mean   median       uq      max neval
##  default_method 365.8609 622.4876 858.0973 703.6735 818.1578 5429.248   100
##     fast_method 133.0536 208.8754 302.1371 248.1200 294.6087 2140.020   100
\end{verbatim}

The fast method is roughly 2.5 times faster than the default method.

\subsection{g.}\label{g.}

Now consider vectorizing the function in point c.~Compare this approach
to the other two. Comment on the results.

\begin{Shaded}
\begin{Highlighting}[]
\NormalTok{run\_kruskal\_vec }\OtherTok{\textless{}{-}} \ControlFlowTok{function}\NormalTok{(X, g) \{}
\NormalTok{  n }\OtherTok{\textless{}{-}} \FunctionTok{ncol}\NormalTok{(X)}
\NormalTok{  m }\OtherTok{\textless{}{-}} \FunctionTok{nrow}\NormalTok{(X)}
  
\NormalTok{  R }\OtherTok{\textless{}{-}} \FunctionTok{t}\NormalTok{(}\FunctionTok{apply}\NormalTok{(X, }\DecValTok{1}\NormalTok{, rank))}
  
\NormalTok{  R\_sum1 }\OtherTok{\textless{}{-}} \FunctionTok{rowSums}\NormalTok{(R[, g }\SpecialCharTok{==} \DecValTok{1}\NormalTok{])}
\NormalTok{  R\_sum2 }\OtherTok{\textless{}{-}} \FunctionTok{rowSums}\NormalTok{(R[, g }\SpecialCharTok{==} \DecValTok{2}\NormalTok{])}
\NormalTok{  R\_sum3 }\OtherTok{\textless{}{-}} \FunctionTok{rowSums}\NormalTok{(R[, g }\SpecialCharTok{==} \DecValTok{3}\NormalTok{])}
  
\NormalTok{  H\_vec }\OtherTok{\textless{}{-}}\NormalTok{ (}\DecValTok{12} \SpecialCharTok{/}\NormalTok{ (n }\SpecialCharTok{*}\NormalTok{ (n }\SpecialCharTok{+} \DecValTok{1}\NormalTok{))) }\SpecialCharTok{*}\NormalTok{ (R\_sum1}\SpecialCharTok{\^{}}\DecValTok{2}\SpecialCharTok{/}\DecValTok{20} \SpecialCharTok{+}\NormalTok{ R\_sum2}\SpecialCharTok{\^{}}\DecValTok{2}\SpecialCharTok{/}\DecValTok{20} \SpecialCharTok{+}\NormalTok{ R\_sum3}\SpecialCharTok{\^{}}\DecValTok{2}\SpecialCharTok{/}\DecValTok{10}\NormalTok{) }\SpecialCharTok{{-}} \DecValTok{3} \SpecialCharTok{*}\NormalTok{ (n }\SpecialCharTok{+} \DecValTok{1}\NormalTok{)}
  
  \FunctionTok{return}\NormalTok{(H\_vec)}
\NormalTok{\}}

\CommentTok{\# benchmark}
\NormalTok{performance\_all }\OtherTok{\textless{}{-}} \FunctionTok{microbenchmark}\NormalTok{(}
  \AttributeTok{default\_method =} \FunctionTok{run\_kruskal\_default}\NormalTok{(X, grp),}
  \AttributeTok{fast\_method    =} \FunctionTok{run\_kruskal\_fast}\NormalTok{(X, grp),}
  \AttributeTok{vec\_method     =} \FunctionTok{run\_kruskal\_vec}\NormalTok{(X, grp),}
  \AttributeTok{times =} \DecValTok{100}
\NormalTok{)}
\FunctionTok{print}\NormalTok{(performance\_all)}
\end{Highlighting}
\end{Shaded}

\begin{verbatim}
## Unit: milliseconds
##            expr      min        lq      mean    median       uq       max neval
##  default_method 398.5373 647.27635 746.20044 717.13890 815.1060 1936.1393   100
##     fast_method 147.4938 219.15465 260.26023 236.22755 286.0756  477.6167   100
##      vec_method   9.9725  20.66925  27.72508  25.49925  30.2897  104.1825   100
\end{verbatim}

The vectorized function is the clear winner by being 21 times faster
than the default method.

\end{document}
